% 《大学物理学学习指导》第5版例16.8—16.10，
% 原书印刷229／PDF物理237页；衰变与粒子反应近似均显式注明。
\subsection{例题：镭活度、电子对产生与湮没}

\sourceexample{16.8}{229}{237}
\problemtext 测得$1$ g镭-226的活度为$1$居里，
$1\,\mathrm{Ci}=3.7\times10^{10}\,\mathrm{s^{-1}}$。
求其半衰期及平均寿命。
\workings
以摩尔质量近似为226g/mol，
$N=N_A(1/226)$，
其中$N_A=6.022\times10^{23}\,\mathrm{mol^{-1}}$。
由$A=\lambda N$，
\[
 \lambda=\frac{A}{N}
 =\frac{3.7\times10^{10}\cdot226}{6.022\times10^{23}}
 \simeq1.39\times10^{-11}\,\mathrm{s^{-1}}.
\]
从而
\[
 \boxed{T_{1/2}=\frac{\ln2}{\lambda}
 \simeq4.99\times10^{10}\,\mathrm s
 \simeq1.58\times10^3\,\mathrm{yr}},
\]
\[
 \boxed{\tau=\frac1\lambda
 =\frac{T_{1/2}}{\ln2}
 \simeq2.28\times10^3\,\mathrm{yr}}.
\]
\crosscheck 活度单位是每秒衰变数而非单核总数。
计算使用原书给定居里换算值，结果与公认的
镭-226约1600年半衰期量级相符。

\sourceexample{16.9}{229}{237}
\problemtext 高能光子与重原子核相互作用，
生成一对电子和正电子。
求能量阈值，并说明原书$1.022$ MeV近似的成立条件。
\workings
电子与正电子静能之和为$2m_ec^2$，
因此在忽略重核反冲的近似下
\[
 \boxed{E_{\gamma,\min}\simeq2m_ec^2
 =1.022\,\mathrm{MeV}}.
\]
若核在实验室原来静止，核质量有限为$M$，
应使用入射的光子与原子核构成的总四动量不变量。
入射不变量为
$s=M^2c^4+2Mc^2E_\gamma$；
阈值时最终质心系粒子静质量总和为$M+2m_e$，
于是
\[
 M^2c^4+2Mc^2E_{\gamma,\min}
 =(M+2m_e)^2c^4,
\]
即
\[
 \boxed{E_{\gamma,\min}
 =2m_ec^2\left(1+\frac{m_e}{M}\right)}.
\]
\crosscheck 原书$1.022$MeV仅在$M\gg m_e$
可忽略核的反冲能量时成立。
一个光子在空旷真空中不能单独转化为正负电子对
而同时满足能量与动量守恒，需要核等第三体。

\sourceexample{16.10}{229}{237}
\problemtext 电子与正电子相遇并湮没，
说明通常产生的辐射及静止电子正电子对的
两光子能量与方向。
\workings
低能、近静止的电子正电子对发生典型
两光子湮没过程：
$e^-+e^+\to\gamma+\gamma$。
总三动量为零时，两光子的动量必须大小相等、
方向相反，且能量守恒给
\[
 \boxed{\bm p_{\gamma1}=-\bm p_{\gamma2}},
 \qquad
 \boxed{E_{\gamma1}=E_{\gamma2}
 =m_ec^2\simeq0.511\,\mathrm{MeV}}.
\]
\crosscheck 原书所说的“相反方向、两束511keV光子”
需要初始正负电子对总动量为零且忽略初始动能。
若初始总动量不为零，实验室系出射角和光子能量
不必相等或严格背对背；不能将近静止情形无限推广。
